\documentclass[11pt,a4paper,oneside]{book}
\usepackage[utf8]{inputenc}
\usepackage[T1]{fontenc}
\usepackage{lmodern}
\usepackage{microtype}
\usepackage[margin=2.6cm]{geometry}
\usepackage{booktabs}
\usepackage{longtable}
\usepackage{array}
\usepackage{xcolor}
\usepackage{listings}
\usepackage{enumitem}
\usepackage{titlesec}
\usepackage{fancyhdr}
\usepackage{textcomp}
\usepackage{amsmath}
\usepackage[hidelinks,pdftitle={nodwyn},pdfauthor={nodwyn project}]{hyperref}

\definecolor{codebg}{gray}{0.96}
\definecolor{codecomment}{rgb}{0.35,0.45,0.35}
\definecolor{codekw}{rgb}{0.15,0.15,0.55}
\definecolor{codestr}{rgb}{0.5,0.2,0.2}

\lstdefinelanguage{Scheme}{
  morekeywords={define,lambda,if,let,let*,letrec,begin,set!,quote,quasiquote,unquote,cond,case,else,and,or,do,when,unless,guard,register-tool,load},
  alsoletter={-!?*>},
  morecomment=[l]{;},
  morestring=[b]",
  sensitive=true
}
\lstset{
  basicstyle=\ttfamily\small,
  backgroundcolor=\color{codebg},
  commentstyle=\color{codecomment}\itshape,
  keywordstyle=\color{codekw}\bfseries,
  stringstyle=\color{codestr},
  columns=fullflexible,
  keepspaces=true,
  breaklines=true,
  frame=none,
  xleftmargin=0.6em,
  framexleftmargin=0.6em,
  aboveskip=0.8em,
  belowskip=0.8em,
  showstringspaces=false,
  tabsize=4,
  extendedchars=true,
  literate={é}{{\'e}}1 {É}{{\'E}}1 {ç}{{\c{c}}}1
}
\lstnewenvironment{go}{\lstset{language=Go}}{}
\lstnewenvironment{scm}{\lstset{language=Scheme}}{}
\lstnewenvironment{sh}{\lstset{language=bash,keywordstyle=\ttfamily}}{}
\lstnewenvironment{txt}{\lstset{language={}}}{}

\newcolumntype{L}[1]{>{\raggedright\arraybackslash}p{#1}}
\newcommand{\nodwyn}{\textsf{nodwyn}}
\newcommand{\code}[1]{\texttt{#1}}
\newcommand{\file}[1]{\texttt{#1}}
\newcommand{\adr}[1]{ADR~#1}
\newcommand{\spec}{\textsc{spec}}
\newcommand{\todo}[1]{\textcolor{red!70!black}{[#1]}}
\newenvironment{aside}{\begin{quote}\small}{\end{quote}}

\titleformat{\chapter}[display]{\normalfont\huge\bfseries}{\chaptertitlename\ \thechapter}{12pt}{\Huge}
\pagestyle{fancy}
\fancyhf{}
\fancyhead[L]{\small\leftmark}
\fancyhead[R]{\small\thepage}
\renewcommand{\headrulewidth}{0.3pt}
\setlength{\parskip}{0.45em}
\setlength{\parindent}{0pt}
\setlength{\emergencystretch}{3em}
\setlist{nosep}

\title{\Huge\textbf{\nodwyn}\\[0.6em]\Large a minimal, self-extending coding-agent harness\\[0.4em]\large the idea, the history, the code, and the language}
\author{the \nodwyn{} project}
\date{September 2026}

\begin{document}
\frontmatter
\maketitle

\chapter*{Preface}
\addcontentsline{toc}{chapter}{Preface}

This book is written for one reader at three moments: the person who wants
to understand what \nodwyn{} is trying to be, the person who has to change its
code, and the person sitting at the prompt. It is organised accordingly.

\textbf{Part I} states the idea: why a coding agent's secrets should be
\emph{credentials} the harness never sees rather than \emph{content} it tries
to recognise; why one guard at one boundary is worth more than many checks
in many places; why the log, not the model's context, is the truth; and
every channel through which bytes enter or leave a session\,---\,models, MCP
servers, keyed APIs, the judge, files, programs, you, and the documents you
attach. \textbf{Part II} is how \nodwyn{} was built: bootstrapped with
another coding-agent harness until it could take over its own development,
then grown phase by phase; and what is still to come. \textbf{Part III} is
an inventory of what exists. \textbf{Part IV} walks the code: one tool call
followed from the keyboard to the disk and back, then one that leaves the
machine. \textbf{Part V} is a tutorial and reference for \nodwyn{} Scheme,
the small language in which \nodwyn's own tools, and its agent loop, are
written. \textbf{Part VI} is the operator's manual. For a slower
introduction aimed at users, see \emph{Becoming a \nodwyn{} Expert}
(\file{docs/expert/}).

Everything in this book is derived from the repository: \file{SPEC.md},
\file{plan.md}, \file{docs/}, \file{docs/adr/}, and the source. Where the
book and the repository disagree, the repository is right and the book has
a bug.

\paragraph{Status.} The code described is the main branch as of
2026-09-30: the last promotion tag, \code{v0.42-console}, and what has been
merged since. Promotions no longer create tags; a stable binary's version is
the commit it was built from.

\tableofcontents

\mainmatter
\part{The idea}

\chapter{What \nibb{} is}

\nibb{} is a coding agent's \emph{harness}: the program that sits between a
language model and your machine. It holds the conversation, offers the model
tools, runs the tools the model asks for, and feeds the results back. Every
such harness answers three questions, and \nibb{} exists because it answers
them differently.

\paragraph{Who decides what the model may do?} In \nibb{}, one function. Every
tool the model can call is a \emph{primitive} in an embedded Scheme
interpreter, and every primitive application, whoever asks for it, passes
through a single Go function where a \emph{guard} decides: allow, ask the
user, or deny. There is no second path. The Scheme interpreter is the only
surface through which anything reaches the host, and the guard sits on it.

\paragraph{How are secrets protected?} By never being content. A key is not
a string the harness tries to spot in a request; it is a \emph{credential}
the user registered into an encrypted store, referenced by an opaque handle,
and injected only at a destination the user has \emph{bound} it to. The
question ``is this egress safe?'' is not ``does this look like a secret?''
but ``is this class of credential authorised for this destination, in this
position, with this method?''\,---\,a table lookup, owned by the user. Content
detection exists too (a cheap deterministic layer, a taint layer, an optional
learned layer), but it is the residual, not the authority.

\paragraph{What is the source of truth?} An append-only log. One JSON line
per event, fsynced. The model's context window, the transcript you see, the
summaries a later phase will add\,---\,all are derived from the log and can
be recomputed from it. A crash loses nothing; a replay reproduces a session
byte for byte.

\paragraph{And one more thing.} \nibb{} extends itself. The tools it uses to
build itself\,---\,run the tests, show a diff, search the tree\,---\,are written in
its own Scheme, loaded from the repository, and cross the same guard as
everything else. From the fourth week of its life, \nibb{} has been the tool
used to build \nibb{}.

\section*{Shape}

\begin{itemize}
\item \textbf{Stack.} Go host; a user-written Scheme interpreter embedded in
  it; a single static binary; a line-mode front-end now, a TUI later, and a
  core that knows nothing about either.
\item \textbf{Model.} A local Qwen served by Ollama on a workstation on the
  tailnet. No API key, no cloud dependency, and\,---\,until credentials exist\,---\,no
  secret anywhere on the build machine.
\item \textbf{Scope.} One developer, three to six months, v0. No sub-agents,
  no GUI, no embeddings, no cross-session memory. Those are decisions, not
  omissions.
\end{itemize}

\chapter{Invariants}

Five statements are load-bearing. Everything else in the design is a
consequence of them, and every phase of the plan is judged by whether it
preserves them.

\begin{enumerate}[leftmargin=2em]
\item \textbf{The log is the single source of truth.} Index and window are
  derived, revocable views; a stale summary means recompute-from-log.
\item \textbf{The judge judges intent and plans over redacted content; the
  local stack judges \emph{what}.} Jev\,---\,the external calibrated judge used
  at two seams\,---\,sees the ask, the plan, and egress \emph{situations}
  (class, destination, scope), after redaction. It never sees a secret
  value. It is never audited for the harness's own control-plane key.
  (Wording amended by \adr{19}.)
\item \textbf{Secrets are credentials, not content.} The harness holds opaque
  handles into a log-excluded encrypted keychain, uses them only at a bound
  destination, and never sees the value.
\item \textbf{The Scheme interpreter is the only Scheme$\to$Go surface; the
  guard sits on it.}
\item \textbf{Every I/O edge is guarded\,---\,by the mechanism that fits its
  substrate.} Scheme values get provenance tags; text gets redaction; egress
  gets the binding check.
\end{enumerate}

Invariant~4 is the one you will meet first in the code: \code{applyPrim} in
\file{internal/scheme/prim.go} is the only caller of any primitive's handler,
and a test records every crossing to prove it. Invariant~1 is the one that
shapes everything you will not see: there is no ``conversation state'' object
anywhere; there are events, and functions that fold them.

\chapter{Secrets are credentials}

\section{Binding, not recognition}

A secret is never suspicious in isolation. The bytes \code{sk-ant-...} sent to
\code{api.anthropic.com} are correct; the same bytes sent to
\code{collect.example.com} are exfiltration. Rightness is a \emph{binding}
question:

\begin{quote}
Pass a secret-bearing egress \emph{iff} its class is authorised for
(destination, positions, methods).
\end{quote}

The binding table is user-owned: rows of \code{(class $\to$ destination,
positions, methods)}. A bound egress passes silently and is logged. An
unbound or unknown one is denied by default (a confirm is opt-in policy)\,---\,it
never fails open. The harness does not \emph{classify} a secret; the user
\emph{binds} one, and the class is a name the user chose when registering it.

\section{Two planes}

The harness's own model key\,---\,when it has one\,---\,is \emph{control plane}:
its identity, its standing; not something it can grant or revoke, and not
something the judge audits. Every \emph{other} credential the user registers
is \emph{data plane}: bound like any other, injected only at its destination.
Under this split, other LLM providers become tools among tools: an OpenAI or
Anthropic call is an egress with a bound credential, judged exactly like an
HTTP request to anywhere else. In the bootstrap, the control plane is a local
model that needs no key at all, so the first real secret on the machine will
be the judge's own API key, in Phase~5.

\section{The content-detection ladder}

Binding is the authority, but three local, zero-retention layers watch
content, in a fixed order, each cheaper than the next and each able to stop
the escalation:

\begin{center}
\begin{tabular}{@{}lp{2.2cm}p{7.2cm}p{2.6cm}@{}}
\toprule
Layer & Acts on & Fires when & Cost \\
\midrule
L1 & text I/O & a known secret \emph{shape}: regex, entropy, prefix. Fails closed. & sub-millisecond \\
L2 & Scheme values & a value tagged at ingress as a known class heads for egress. O(1) flag at the FFI. & nothing \\
L3 & text I/O & only the class-unknown residual L1/L2 did not decide. A calibrated $P(\text{secret})$ that feeds the decision with binding, never alone. & learned; optional \\
\bottomrule
\end{tabular}
\end{center}

The deterministic layers early-terminate the learned one: if L1 or L2 already
\emph{knows}, L3 is not paid for. L2 is scoped honestly\,---\,it stops the
\emph{model's} indirect leak, not a user who can strip a tag on purpose, and
the spec rates it accordingly (0.54).

\section{The four edges}

Every path between the outside world and the two substrates\,---\,Scheme values
and the model's text context\,---\,is an edge with its own mechanism:

\begin{center}
\begin{tabular}{@{}lllp{6cm}@{}}
\toprule
Edge & From $\to$ to & Substrate & Mechanism \\
\midrule
E1 & external $\to$ Scheme value & scheme & L2 taint (provenance tag) \\
E2 & external $\to$ model context & text & L1+L3 redaction to \code{[REDACTED class]} \\
E3 & model/Scheme $\to$ external & both & binding check + L1/L3 (+ L2 flag) \\
E4 & user/OS $\to$ model context & text & L1+L3 redaction (catches pasted secrets) \\
\bottomrule
\end{tabular}
\end{center}

Redaction is a \emph{substrate controller}: it runs at every ingress to the
context, not just on API responses, so a leaked value can never be
re-emitted, even to a bound destination. The category error the design had
to resolve was thinking a taint on a value could protect the context window;
it cannot (a window is text, not a value), so the text layers run there.

\section{The judge: where it is allowed}

Jev is an external calibrated judge with no retention guarantee. The
original spec drew the line at ``never content''; \adr{19} redrew it where
it belongs, at ``never a secret value''. The question at Seam A\,---\,\emph{is
this plan good enough for the user's ask?}\,---\,has no answer without the ask
and the plan, and both are content.

So Jev is treated as what it is: an \emph{egress edge}. Its credential is a
bound data-plane credential (E3), and its payload passes the same
redaction controller as the model context (E2/E4) before it leaves. No
bespoke sanitiser; two existing mechanisms composed. What may go beyond
the ask and the plan is a user policy\,---\,\code{jev.content}: \code{none},
\code{plan} (default), \code{plan+summaries}, or \code{full}, the last
opt-in per repository because it is the level at which proprietary code is
retained elsewhere\,---\,and every call is logged with its level, size, payload
hash and the redactions it triggered.

Two seams:

\begin{description}[leftmargin=1.6em]
\item[Seam A] score a draft \emph{plan} against the \emph{ask} for quality,
  coherence and fulfilment; the model revises on a low score, at most three
  rounds.
\item[Seam C] score an ambiguous egress \emph{situation}\,---\,class,
  destination, scope\,---\,never the value.
\end{description}

Never: a secret value (enforced by redaction, not by withholding content),
control flow, its own loop, or the harness's control-plane credential. The
principle, restated: the judge judges intent and plans over \emph{redacted}
content; the local stack judges provenance and secret shape, and is the
only thing that ever sees a value.

\chapter{The loop and the guard}

The agent loop is the plain one: ask the model; run the tool it asks for,
through the guard; feed the result back; ask again. Stop on a plain answer or
a hard stop. No sub-agents in the core loop.

The loop is itself a Scheme procedure, \code{agent-turn}, so how a turn goes
can change without changing Go: plan first, check with the judge, stop a
batch of calls at the first failure. What a turn must not skip stays in Go,
behind a few primitives: windows built from the log, only the model's own
calls run, exactly one result per call, and budgets and a progress guard
that the procedure cannot catch (\adr{26}). The same procedure can route
work: narrow the context to a task and a plan, and hand it to another model
the user configured (\adr{27}).

Before each message, a plan gate: a yes/no from the judge on the abstract
intent (``is this multi-step?''), logged per message, plus a code-enforced
hard trigger for irreversible or high-stakes actions (force pushes,
recursive deletes, writes outside the tree, any unbound egress). The gate is
chosen over a pure pull model because it is auditable per message and
improves as the judge improves.

The guard is a pipeline of \emph{checkers}. Each looks at the primitive's
declared metadata and its actual arguments and returns a verdict; the
strictest wins (deny over confirm over allow). A confirm is put to the user
through a \emph{confirmer} the core implements as an event plus a wait, so
the guard never knows what a terminal is. A \emph{standing approval}\,---\,a
command prefix or a write path the user authorised ahead of time\,---\,turns a
confirm into an allow and never turns a deny into anything. Every decision
that reaches the host, and every non-allow, is logged with who asked (the
user, the model or the harness) and why.

\chapter{Memory}

Three tiers, one truth:

\begin{itemize}
\item \textbf{Disk}: the lossless log plus a session-scoped structural index.
  Source of truth.
\item \textbf{Cache}: summaries, written \emph{as log entries}, rebuildable and
  revocable, possibly stale.
\item \textbf{RAM}: the token-budgeted window, with surgical pull on demand.
\end{itemize}

v0 has no cross-session memory and no embedding index; retrieval is text
search (grep). The bootstrap ships a deliberately naive
window\,---\,elide the oldest tool results, then the oldest turns\,---\,but even it is
rebuilt from the log on every request, and it elides \emph{in place} so the
request prefix never changes and the model server's cache keeps hitting.

\chapter{Honest limits}

The spec records its own confidence, and the low numbers are the useful
ones.

\begin{enumerate}[leftmargin=2em]
\item L2 is scoped, not a security boundary. It stops the model's
  unintended leak; a user can strip a tag. The real lines are binding
  (E3), fail-closed defaults, and ingress redaction (E2/E4).
\item A bare ``confirm'' can be uninformed consent. If the class is also
  unknown, the prompt is a rubber stamp. Mitigation: always name class and
  destination; prefer requiring a binding to a bare yes.
\item L3's value is its training data: strong for known shapes, weak for a
  novel internal token\,---\,exactly the residue binding catches.
\item ``The egress surface is closed'' stands at 0.30 until redaction is a
  substrate controller and E4 is covered. Once all four edges are guarded,
  the honest statement is: no path exists for a secret \emph{value} to egress
  unless bound or explicitly re-confirmed\,---\,but ``closed'' is the user's
  affirmation, not something the harness can prove of itself.
\end{enumerate}

And in the bootstrap, before any of P3--P7 exist, there is a fifth: the
\code{shell-exec} hole. A \code{y} to \code{cat \textasciitilde/.ssh/id\_rsa}
is a \code{y}. The prompt names the path; the operator reads it.

\part{Where things are going}

\chapter{Bootstrap, then self-host}

The plan has two stages. \textbf{Stage 0} (weeks 1--4) is hand-written: the
smallest \nibb{} that can do real work on its own repository. \textbf{Stage
1+} (weeks 5--24) is every remaining phase of the spec's build order, executed
as \nibb{} sessions driving a promoted binary. The spec's ordering is kept;
only the entry point changed.

\section{Why this order}

Three deliberate deviations from the spec's numbered build order:

\begin{itemize}
\item \textbf{The log first, not last.} Invariant~1 makes it the source of
  truth and every later phase says ``logged''. A 200-line JSONL writer on day
  one beats retrofitting six phases.
\item \textbf{A walking-skeleton loop right after the guard.} The security
  phases are then tested against a real loop, not a mock. The TUI stays at
  its spec position as a second front-end over the same events.
\item \textbf{A minimal kernel, grown later by \nibb{} itself.} To self-host
  you need a working loop before a complete interpreter. Growing the
  language is the ideal first dogfood task: pure Go, crisp tests.
\end{itemize}

\section{Two binaries, always}

\file{bin/nibb} is the working-tree build. \file{\textasciitilde/.local/bin/nibb-stable}
is the last \emph{promoted} build. All development runs \code{nibb-stable};
\nibb{} never edits itself through the binary being edited, so a broken tree
can never strand you. Promotion (\code{make promote TAG=v0.N-phase})
refuses a dirty tree, runs the acceptance suite and a self-check\,---\,the
candidate completes a scripted session and replays its own log\,---\,then swaps
the binaries, keeps the previous one as \code{.prev}, and tags.

\section{The safety envelope}

Using \nibb{} before the keychain, L1 and redaction exist requires
compensating controls. Each is retired by the phase that builds the real
thing; some stay as defense in depth.

\begin{center}
\small
\begin{tabular}{@{}p{6.4cm}p{3.4cm}p{4.6cm}@{}}
\toprule
Bootstrap control & Stands in for & Retired by \\
\midrule
No credential of any kind on the build machine; the control plane needs none & keychain, bindings & P3 builds the keychain; P5 registers the first real credential \\
Child processes get an allowlisted environment only & E2 redaction of \code{env} & never\,---\,defense in depth \\
No \code{http-request} primitive & binding table / E3 & P3 \\
Paths pinned to the repository root, symlinks resolved & --- & never (per-project cwd is the model) \\
Path denylist: \code{.env*}, \code{*.pem}, \code{*.key}, \code{id\_*}, \code{\textasciitilde/.ssh}, \code{\textasciitilde/.aws}, \ldots & L1 / E2 & never\,---\,defense in depth \\
Command-line trigger: \code{curl wget nc ssh scp rsync}, \code{git push}, \code{rm -rf}, absolute paths outside the repo $\Rightarrow$ confirm & L1 on argv, P5 hard triggers & P5 \\
Irreversible primitives $\Rightarrow$ confirm, with a diff for writes & --- & permanent \\
Operator rule: no secrets in the repo; never paste one into \nibb{} & E4 redaction & P7 \\
Operator rule: bootstrap-era \nibb{} only on trusted repositories & everything & P7 \\
\bottomrule
\end{tabular}
\end{center}

\section{Milestones}

\begin{center}
\begin{tabular}{@{}llll@{}}
\toprule
Weeks & Phase & Promotes & Milestone \\
\midrule
1--4 & Bootstrap B0--B3 (hand-written) & \code{v0.1-bootstrap} & \textbf{M0 self-hosting} (done 2026-09-21) \\
5--6 & P1b kernel growth & \code{v0.2-kernel} & committed 2026-09-21; promotion pending \\
7--9 & P3 credentials & \code{v0.3-creds} & \\
10--11 & P4 L1 & \code{v0.4-l1} & \textbf{M2 safe with real keys} \\
12--13 & P5 plan gate + Jev & \code{v0.5-plan} & \\
14--15 & P6 L2 taint & \code{v0.6-taint} & \\
16--17 & P7 redaction & \code{v0.7-core} & \textbf{M3 v0-core} (3--4 month cut) \\
18--19 & P8 L3 (optional) & \code{v0.8-l3} & \\
20--22 & P9 TUI & \code{v0.9-tui} & \\
23--24 & P10 memory & \code{v0.10} & \textbf{v0} \\
\bottomrule
\end{tabular}
\end{center}

If time runs short, cut in this order: P10, P8, P9 polish. Never P3--P7.

\section{The friction log}

When \code{nibb-stable} gets in the way, the rule is: hand-edit, and write
one line in \file{docs/friction.md}. That file is the backlog for the
front-ends and for memory, and it already paid for itself: the first
self-hosting run showed the model rewriting a 400-line file to change one
line; the entry became \code{fs-edit} the same afternoon, and the second run
took a third of the tool calls. Fallback is allowed and honest: the goal is a
built \nibb{}, not purity.

\chapter{The phases}

Each phase is a set of \nibb{} sessions plus a promotion. The ``with \nibb{}''
note says what the tool can do at that point and what its promotion unlocks.

\paragraph{P1b\,---\,grow the kernel (weeks 5--6).} \code{letrec}, \code{case},
\code{do}, \code{when}/\code{unless}, vectors, a string port, \code{guard} for
catching errors, block comments. Done 2026-09-21 by an external agent as a
hand-coded fallback (\adr{16}), then reviewed: the review found a crash
(\code{(when)}), a \code{guard} that could swallow denials and aborts, and a
nonstandard \code{guard} shape; all fixed before commit. The conformance
suite roughly doubled.

\paragraph{P3\,---\,credentials (weeks 7--9).} The encrypted keychain
(XChaCha20-Poly1305, argon2id), \code{credential-register} with masked input
returning an opaque handle, the binding table, the E3 check in the guard,
\code{http-request} injecting handles only at bound positions, and a second
provider for data-plane \code{llm-call}. Proven by the first real exfil
scenarios: bound header passes; other host, or body instead of header,
denied; the value appears nowhere.

\paragraph{P4\,---\,L1 (weeks 10--11).} The shape catalog (prefixes, entropy,
structure), the egress allowlist, a checker over text arguments. Sub-millisecond
at p99 on 64\,KB; a false-positive sweep over real repositories. \nibb{} runs
the sweep on itself.

\paragraph{P5\,---\,plan gate and Jev (weeks 12--13).} The System One client,
the noul gate, Seam A scoring with revision rounds, the hard-trigger list.
Jev is an egress edge (\adr{19}): its credential is bound, its payload\,---\,the
ask and the plan, by default\,---\,is redacted by the same controller the
model context gets, and every call is logged with what left. The Jev key is
the first real secret on the machine. After promotion, every later build
session is planned.

\paragraph{P6\,---\,L2 taint (weeks 14--15).} The reserved \code{Taint} field is
populated at E1, propagated once in the dispatch path, checked at E3; when
L2 knows, L3 is not invoked. \code{box}/\code{wipe} strip a tag only from
user-originated evaluation. This phase edits the interpreter's own dispatch\,---\,safe
only because of the two-binary rule.

\paragraph{P7\,---\,ingress redaction (weeks 16--17).} One context-ingress
function, enforced by lint; \code{[REDACTED class]} at E2 and E4; the same
redactor in front of the log writer. This retires the last operator rules.
\nibb{} is now safe on other repositories with real keys, and
\textbf{M3 v0-core} is reached.

\paragraph{P8\,---\,L3 (weeks 18--19, optional).} Features to logistic regression
with calibration; trained on a public corpus plus \nibb's own confirmed
events; invoked only on the residual; excluded with a build tag.

\paragraph{P9\,---\,TUI (weeks 20--22).} A bubbletea front-end over the same
\code{core.Session}, with no core change; the friction log is its backlog.

\paragraph{P10\,---\,memory (weeks 23--24).} The spec's window: summaries as
revocable log events, a tree-sitter and grep index exposed as primitives,
byte-for-byte replay of the window. Designed against real logs of \nibb{}
building \nibb{}.

\chapter{The control plane}

The model is \code{qwen3.8-27b-96k} served by Ollama~0.34 on \code{workshop},
reachable only inside the tailnet, with a 98,304-token context, one loaded
model and one request slot. Verified facts and their consequences:

\begin{itemize}
\item \textbf{One slot}, shared with other clients. The loop tolerates queueing:
  long per-request timeouts, streaming keeps the front-end alive, abort
  cancels the HTTP request so the slot frees.
\item \textbf{Prefix stability is a performance feature.} A cold prompt
  evaluation of a large prefix costs tens of seconds; a warm one is free.
  So: system prompt, then tool schemas in stable order, then messages in
  append-only order; elide in place; load tools before the first message.
  The native API reports evaluation time, not cached tokens; sub-second
  \code{prompt\_eval\_ms} is the cache signal.
\item \textbf{Thinking is data, not context.} Logged, shown dimmed, never fed
  back.
\item \textbf{A 27B model is the dogfood ceiling.} It does well with a
  paragraph and badly with a page; tasks are small and explicit; the plan
  gate does more work here than it would for a frontier model.
\item \textbf{Swapping the control plane is a config change} once P3 exists\,---\,a
  decision to be made on friction-log evidence, not up front.
\end{itemize}

\chapter{Long-running turns}

Today a turn stops after fifty tool calls. That is the right \emph{kind} of
guard and the wrong \emph{unit}. \adr{18} (proposed) replaces it with:

\begin{itemize}
\item \textbf{Budgets, not a count}: tool calls, wall clock, model tokens and
  irreversible actions, each with a soft limit that raises a
  \emph{checkpoint} and a hard limit that stops.
\item \textbf{Checkpoints through the confirmer}: the same channel as a guard
  prompt, so any front-end can answer and an unattended run gets a policy
  answer; each is a log event.
\item \textbf{A progress guard}: the same call four times in a row, a denied
  call retried, ten consecutive tool errors\,---\,stop early, say why.
\item \textbf{Long runs are chained turns}: the loop stays one turn; a driver
  in the front-end sends ``continue'' while the run budget allows. Because
  the log is the truth, a run stops and resumes at any turn boundary.
\item \textbf{Unattended means more guards}: no interactive confirms (so every
  write needs a standing approval), the progress guard on, hard budgets, a
  plan whose steps are checkpoint boundaries, and the same envelope.
\end{itemize}

\part{What exists today}

\chapter{Inventory}

\section{Packages}

{\small
\begin{longtable}{@{}L{4.8cm}L{10.2cm}@{}}
\toprule
Package & Owns \\
\midrule
\endhead
\code{internal/scheme} & the language: values and their provenance tags, reader, \code{eval}/\code{apply}, builtins, the primitive table, \code{applyPrim} \\
\code{internal/guard} & the decision pipeline and its checkers: metadata, the pin and mounts, the path denylist, command-line triggers, standing approvals, E3, L1, L2, user checkers \\
\code{internal/prim} & host primitives: files, search, git, programs, documents, \code{http-request}, the judge, policy files, the version store \\
\code{internal/keychain}, \code{internal/binding} & encrypted credentials and their injection; the binding table \\
\code{internal/detect/l1}, \code{internal/redact} & credential-shape detection; redaction of text \\
\code{internal/llm} & the \code{Provider} interface; Ollama and chat-completions clients; what a request cost \\
\code{internal/mcp} & a client for MCP servers, over streamable HTTP or a contained server's standard input and output; manifest hashing \\
\code{internal/container}, \code{internal/egress} & running a contained MCP server with podman or docker; the egress proxy that is its only way out \\
\code{internal/extract} & text from attached documents; checks on attached images; every text a document carries, before it is uploaded \\
\code{internal/window} & the context window, derived from the log \\
\code{internal/loop} & the loop's primitives and its default Scheme: the strategy, the plan gate, prompts, compaction, sub-agents, working habits \\
\code{internal/budget} & per-turn accounting, ceilings and checkpoints \\
\code{internal/core} & \code{Session}: the UI-agnostic fa\c{c}ade that assembles the rest \\
\code{internal/event}, \code{internal/log} & what a front-end sees; the append-only log and its typed payloads \\
\code{internal/config}, \code{internal/policyfile} & paths and settings files; locked, atomic policy-file writes \\
\code{internal/frontend/plain} & line mode, the Scheme REPL, replay \\
\code{internal/frontend/tui} & the full-screen terminal client \\
\code{internal/frontend/rpc} & the daemon, its protocol and sessions \\
\code{internal/frontend/web} & the browser console and its bridge \\
\code{internal/frontend/run} & the driver of \code{nodwyn run}: a task carried to its end with no conversation, for scripts and benchmarks \\
\code{cmd/nodwyn} & flags $\to$ core $\to$ front-end; nothing else \\
\code{cmd/nodwyn-egress} & the egress proxy, built static for Linux to run in a container \\
\code{dev/}, \code{.nodwyn/} & this repository's own tools and standing approvals \\
\code{test/exfil} & the exfiltration scenario suite \\
\code{test/fakes} & scriptable fakes of a model server, the judge, an MCP server, a container runtime, an OAuth2 provider and Google Drive \\
\code{test/browser} & browser flows against the console \\
\bottomrule
\end{longtable}}

\section{Commands}

\begin{txt}
nodwyn                        chat in the current directory (line mode)
nodwyn --repo DIR             pin to another directory
nodwyn --resume ID            continue a session in its recorded directory
nodwyn tui                    full-screen terminal client (local or daemon)
nodwyn daemon                 serve sessions over a private Unix socket
nodwyn ui                     the browser console for a running daemon
nodwyn ui --listen ADDR       also for other devices, paired by invitation
nodwyn app                    daemon, console and window in one program
nodwyn devices [revoke NAME]  list or revoke paired devices
nodwyn devices invite         an invitation for another device
nodwyn run ...                a batch run, for scripts and benchmarks
nodwyn sessions               list sessions
nodwyn replay ID              print a session log as a transcript
nodwyn replay ID --window SEQ rebuild and verify one model request
nodwyn summary ... ID         add, list or revoke a summary
nodwyn versions FILE          a file's kept previous versions
nodwyn undo FILE              put back the newest kept version
nodwyn keychain init          create the encrypted keychain
nodwyn credential ...         list, add (OAuth2), log in again, revoke
nodwyn mcp [pull NAME]        MCP servers; pull a contained one's image
nodwyn repl | eval            Scheme, outside any session
nodwyn paths                  config and data locations
\end{txt}

\section{Tools the model sees}

Every session offers \code{fs-read}, \code{fs-write}, \code{fs-edit},
\code{fs-list}, \code{fs-find}, \code{index-grep}, \code{index-search},
\code{index-update}, \code{recall}, \code{shell-exec},
\code{shell-start}, \code{shell-read}, \code{shell-send},
\code{shell-stop} and \code{doc-export}; in a git work tree also
\code{git-status} and \code{git-diff}; with a stored \code{google}
credential, \code{google-doc-upload}; with the judge's key stored and
\code{judge.content} set to \code{full}, \code{judge-yes-no}, \code{judge-score} and
\code{judge-choice}. The default loop adds four tools of its own, written
in Scheme: \code{read-definition}, \code{second-look}, \code{plan-write}
and \code{plan-step}. A repository adds
more in its reviewed \file{.nodwyn/init.scm}; after \code{nodwyn trust}, this
one loads \file{dev/}, which offers
\code{go-test}, \code{go-build}, \code{go-vet}, \code{go-fmt-check},
\code{run-make}, \code{git-log}, \code{grep-repo}, \code{find-files} and
\code{read-adr}. Each MCP tool you turned on, on a server whose manifest you
approved and that is offered in this workspace, appears as
\code{server\_\_tool}. \code{:tools} lists what a session offers.

\begin{sloppypar}
Hidden from the model, callable from Scheme: \code{http-request},
\code{credential}, \code{credential-register}, the \code{judge}
primitives below \code{full},
the session's own (\code{scratch-dir}, \code{repository},
\code{workspace}, \code{project-notes}, \code{project-notes-save},
\code{project-instructions}), the
loop's primitives, the policy-file primitives, \code{load},
\code{register-check} and \code{wipe}.
\end{sloppypar}

\section{Guarantees, and the tests that hold them}

{\small
\begin{longtable}{@{}L{5.2cm}L{9.8cm}@{}}
\toprule
Claim & Held by \\
\midrule
\endhead
The log is the single source of truth & \code{nodwyn replay} reproduces a transcript from the file alone; the window takes only events; \code{log\_test.go} round-trips and repairs a truncated tail; \code{replay --window} verifies a request against its recorded digest \\
One Scheme$\to$Go exit, the guard on it & \code{applyPrim} is the only caller of \code{Prim.Handler}; \code{TestGateIsTheSingleExit} records every crossing, including through \code{map}, \code{apply}, \code{sort} and a registered tool \\
MCP tools cross the guard & \code{TestMCPToolsCrossTheGuard}: an undeclared tool confirms with no standing answer, the key reaches the server and neither the log nor the model, the checks are recorded; \code{TestMCPManifestChangeOffersNothing}; for a contained server, \code{TestContainedServerCrossesTheGuard} and \code{TestContainedCredentialMustBeBoundForEveryHost}, against a fake runtime \\
Loading \code{.scm} cannot change evaluation & \code{TestKernelFixed}: only the table grew, \code{define-syntax} stays a binding \\
A strategy cannot skip the invariants & \code{strategy\_test.go}: fabricated or rewritten calls refused, unanswered calls answered, stops uncatchable, provenance kept \\
Proper tail calls & \code{TestTailCalls}: 200k iterations in nine tail positions under an 8\,MB stack cap \\
Every primitive is documented & \code{TestDocCoverage}, \code{TestSessionPrimitivesAreDocumented} \\
A handle never becomes text & \code{TestHandlePrintsOpaque} \\
What enters the context is redacted & \code{TestIngressCanary} and the redaction tests in \file{internal/core} \\
The core knows nothing about terminals & \file{scripts/lint-boundaries.sh} in \code{make lint} \\
No path for the canary to leak & \code{test/exfil}: nineteen scenarios, run with writes in the folder allowed, as sessions have them by default (a write to a hook, a credential-shaped write and a write through a link still ask or are refused), plus Go tests for bindings, redirects, the judge, OAuth2 credentials and document uploads; the scenarios run in a real session, so the guard is the one core assembles; each asserts the decision (and whether a confirm is a security confirm no standing approval may waive) and greps the canary out of the log, the output and the requests to the model \\
Request prefix is byte-stable & \code{TestTurnWithToolCall} compares consecutive requests \\
\bottomrule
\end{longtable}}

\chapter{What is not there yet}

\begin{itemize}
\item No L3. Known key shapes are detected; unknown ones are not.
\item \nodwyn{} runs a local MCP server only in a container, with podman or
  docker, never as a process of yours. A contained server that dies is not
  restarted: its calls fail and say why. One that ignores the proxy
  variables reaches nothing. \nodwyn{} does not
  discover a server's own authorization or register itself as a client: a
  server behind a login takes an OAuth2 credential you add with
  \code{nodwyn credential add -{}-oauth2}.
\item No PDF output built in: \code{doc-export} writes \file{.docx},
  \file{.odt} and \file{.html}, and a PDF is built from LaTeX with
  \code{shell-exec} (\code{latexmk}), since no sandbox covers a TeX engine.
\item Nothing reads an image's pixels: an attached image goes to a model
  unredacted, after a confirm that says so.
\item No embedding index and no structural outline; retrieval is text
  search over the tree, and word search over an index of it, over the
  conversation and, when you allow it for a folder, over its earlier tasks
  (\code{index-search}, \code{index-update}, \code{recall}).
\item One person. A workspace is not shared, and the socket and the paired
  devices are one user's.
\item A spend limit is off unless you set one (\code{cost\_soft} and
  \code{cost\_hard} in \code{[budget]}, or \code{-{}-budget-cost}), and it
  counts only the costs it knows: a request that neither the provider nor
  the model's \code{price\_*} prices adds nothing to it.
\item A plan step marked done is the model's word. An unattended run
  (\code{nodwyn run -{}-unattended}) goes on from turn to turn on those
  marks, and nothing checks them against the transcript. A \code{Test:}
  step must say why its test fails on today's code, but nothing runs the
  test there to see that it does.
\end{itemize}

\part{Reading the code}

\chapter{The spine: one tool call, end to end}

The fastest way into \nodwyn{} is to follow a single tool call from the
keyboard to the disk and back. Eleven files, in order. The excerpts below
are the real code, lightly abridged; line numbers refer to \code{v0.1}.

\section{The front-end sends a line}

\file{internal/frontend/plain/chat.go}. The plain front-end reads a line,
hands it to the session, and then does nothing but render events until the
session says it is idle. It never touches the loop, the guard or the
interpreter\,---\,it is a log subscriber with a reply channel.

\begin{go}
if err := c.Session.Send(line); err != nil {
    fmt.Fprintln(c.Out, "error:", err)
    continue
}
c.follow(ctx)

// follow renders events until the turn ends.
func (c *Chat) follow(ctx context.Context) {
    for {
        select {
        case <-c.Interrupt:
            fmt.Fprintln(c.Out, "\n[aborting]")
            c.Session.Abort()
        case ev, ok := <-c.Session.Events():
            if !ok {
                return
            }
            c.show(ev)
            if ev.Kind == event.Status &&
                ev.Payload.(event.StatusInfo).State == "idle" {
                return
            }
        }
    }
}
\end{go}

\section{The session starts a turn}

\file{internal/core/session.go}. \code{Send} refuses if a turn is running,
otherwise starts one in a goroutine with a cancellable context. When the
turn ends\,---\,however it ends\,---\,the session emits \code{status idle}. Note
\code{turnTrust} being reset: a ``trust writes this turn'' answer cannot
outlive the turn.

\begin{go}
func (s *Session) Send(text string) error {
    s.mu.Lock()
    if s.busy {
        s.mu.Unlock()
        return errors.New("core: a turn is already running")
    }
    ctx, cancel := context.WithCancel(context.Background())
    s.busy, s.turnCancel = true, cancel
    s.mu.Unlock()

    go func() {
        defer func() {
            cancel()
            s.mu.Lock()
            s.busy, s.turnCancel, s.turnTrust = false, nil, false
            s.mu.Unlock()
            s.Emit(event.Event{Kind: event.Status,
                Payload: event.StatusInfo{State: "idle"}})
        }()
        if err := s.loop.Turn(ctx, text); err != nil {
            s.Emit(event.Event{Kind: event.Error,
                Payload: event.ErrorInfo{Err: err.Error()}})
        }
    }()
    return nil
}
\end{go}

\code{Record} is the one function that writes the log, keeps the event for
the window and forwards it to the front-end. It is what the loop, the guard
and the session itself all call:

\begin{go}
func (s *Session) Record(kind log.Kind, payload any) (log.Event, error) {
    ev, err := s.logw.Append(kind, payload)
    if err != nil {
        return ev, err
    }
    s.evmu.Lock()
    s.events = append(s.events, ev)
    s.evmu.Unlock()
    s.Emit(event.FromLog(ev, payload))
    return ev, nil
}
\end{go}

\section{The loop asks the model}

\file{internal/loop/loop.go}. The loop itself is Scheme: \code{Turn} applies
\code{agent-turn}, the default from \file{internal/loop/strategy.scm} or a
project's redefinition, with harness origin, then answers any call the
strategy left unanswered and closes the turn's budget.

\begin{go}
func (l *Loop) Turn(ctx context.Context, request string) error {
    strategy, _ := l.interp.Global().Lookup("agent-turn")
    l.turn = &turnState{budget: budget.From(ctx)}
    defer func() { l.turn = nil }()
    _, err := l.interp.Apply(scheme.WithOrigin(ctx, scheme.OriginHarness),
        strategy, []scheme.Value{scheme.Str(request)})
    if skipErr := l.skipPending("the turn ended before this call ran"); skipErr != nil && err == nil {
        err = skipErr
    }
    if err != nil {
        return surfaced(ctx, err)
    }
    return l.turn.budget.Finish(ctx)
}
\end{go}

The default strategy asks, runs what the model asked for, and asks again
until a reply has no tool calls:

\begin{scm}
(define (agent-turn request)
  (let loop ((reply (model-turn)))
    (let ((calls (assoc-ref reply 'tool-calls '())))
      (unless (null? calls)
        (run-tools calls)
        (maybe-compact request reply)
        (loop (model-turn))))))
\end{scm}

\code{model-turn} (\file{internal/loop/model.go}) is where the model is
asked. Every request rebuilds the window from the events so far (there is no
``messages'' field kept between requests), records an \code{llm.req}
describing the window, streams the response, and records the full
\code{llm.resp}. The budget, the tool-call limit and the progress guard sit
in the primitives, where no strategy can skip them.

\section{The window is a fold over the log}

\file{internal/window/window.go}. \code{Build} turns \code{user.msg},
\code{llm.resp} and \code{tool.result} events into messages and applies a
policy shaped for a prefix cache on one shared model slot: the system
prompt carries the \emph{current task}, so the ask is never trimmed;
elision is \emph{sticky}\,---\,every earlier request's elided set, read back
from the log's \code{llm.req} events, stays elided; and \emph{hysteresis}
trims above a high-water mark down to a low-water mark, oldest first,
tool results before assistant text, user messages never, then holds until
the next crossing. Between trims every request is a byte-identical prefix
plus a delta.

\begin{go}
if total() > pol.HighWater {
    trimmed = true
    protect := len(entries) - Keep
    // Pass 1: oldest tool results. Pass 2: oldest assistant turns.
    for _, kind := range []string{"tool", "assistant"} {
        for i := 0; i < protect && total() > pol.LowWater; i++ {
            if entries[i].kind == kind && !isElided(i) {
                elide(i)
            }
        }
    }
}
\end{go}

\section{The provider streams}

\file{internal/llm/ollama.go}. The native \code{/api/chat} protocol sends
newline-delimited JSON chunks. Content and thinking are accumulated and also
forwarded as deltas; tool calls arrive as JSON objects and are kept as such;
the final chunk carries usage. Nothing in the loop knows this is Ollama.

\begin{go}
for sc.Scan() {
    var ch wireChunk
    if err := json.Unmarshal(sc.Bytes(), &ch); err != nil {
        return nil, fmt.Errorf("ollama: bad chunk: %w", err)
    }
    if ch.Message.Content != "" || ch.Message.Thinking != "" {
        content.WriteString(ch.Message.Content)
        thinking.WriteString(ch.Message.Thinking)
        onDelta(Delta{Text: ch.Message.Content, Thinking: ch.Message.Thinking})
    }
    for _, wt := range ch.Message.ToolCalls {
        out.Message.ToolCalls = append(out.Message.ToolCalls,
            ToolCall{ID: wt.ID, Name: wt.Function.Name, Args: wt.Function.Arguments})
    }
    if ch.Done {
        out.DoneReason = ch.DoneReason
        out.Usage = Usage{PromptTokens: ch.PromptEvalCount, /* ... */}
        break
    }
}
\end{go}

\section{A tool call becomes a Scheme application}

\file{internal/loop/loop.go}, \code{apply}. This is \adr{3} in four lines: look
the primitive up, order the model's named JSON arguments into the
primitive's positional parameters (unknown or missing names are errors the
model is told about), tag the context with \emph{model origin}, and apply.
The model's call now takes exactly the path a call from the REPL would.

\begin{go}
func (l *Loop) apply(ctx context.Context, tc llm.ToolCall) (string, error) {
    p, ok := l.interp.Lookup(tc.Name)
    if !ok || !p.Tool() {
        return "", fmt.Errorf("unknown tool %q; available: %s", tc.Name, l.toolNames())
    }
    args, err := p.ArgsFromJSON(tc.Args)
    if err != nil {
        return "", err
    }
    v, err := l.interp.Apply(scheme.WithOrigin(ctx, scheme.OriginModel),
        scheme.Value{Kind: scheme.KPrim, Prim: p}, args)
    if err != nil {
        return "", err
    }
    return render(v), nil
}
\end{go}

\section{The single exit}

\file{internal/scheme/prim.go}. Whatever the caller\,---\,the loop, \code{map},
\code{apply}, a registered tool's body, the REPL\,---\,every primitive
application reaches this function and no other. Arity first, then the gate,
then the handler. A lint asserts that \code{p.Handler} is invoked nowhere
else.

\begin{go}
// applyPrim is the single Scheme->Go exit. Every primitive application --
// from the model, the prelude, a .scm file or the REPL -- passes through
// here: arity check, guard, handler. Nothing else calls p.Handler.
func (in *Interp) applyPrim(ctx context.Context, p *Prim, args []Value) (Value, error) {
    if err := ctx.Err(); err != nil {
        return Value{}, wrap(p.Name, err)
    }
    if err := checkArity(p, args); err != nil {
        return Value{}, err
    }
    if in.gate != nil {
        if err := in.gate.Before(ctx, p, args); err != nil {
            return Value{}, wrap(p.Name, err)
        }
    }
    v, err := p.Handler(&Call{Ctx: ctx, In: in, Name: p.Name, Args: args,
        Origin: OriginFrom(ctx)})
    if err != nil {
        return Value{}, wrap(p.Name, err)
    }
    return v, nil
}
\end{go}

\section{The guard decides}

\file{internal/guard/guard.go}. Checkers run in order; a deny short-circuits;
the strictest verdict wins; a standing approval converts a confirm into an
allow (and carries the diff into the log, since nobody saw it). A confirm
goes to the confirmer, which the core implements as an event plus a wait.

\begin{go}
for _, c := range g.checkers {
    d := c.Check(ctx, p, args)
    if d.Verdict == Allow {
        if d.Standing && standing == "" {
            standing = d.Reason
        }
        continue
    }
    reasons = append(reasons, d.Reason)
    if d.Detail != "" {
        details = append(details, d.Detail)
    }
    if d.Verdict > final.Verdict {
        final, checker = d, c.Name()
    }
    if d.Verdict == Deny {
        break // nothing can relax a deny
    }
}
if final.Verdict == Confirm && standing != "" {
    final = Decision{Verdict: Allow, Standing: true, Reason: standing}
    reasons, checker = []string{standing}, "standing"
}
\end{go}

\section{A checker reads metadata, not names}

\file{internal/guard/checkers.go}. \code{RepoPin} does not know what
\code{fs-read} is. It asks the primitive's declaration which parameters
carry a path, resolves each against the root\,---\,following symlinks on the
longest existing prefix, so a link out of the tree does not escape\,---\,and
denies anything that lands outside.

\begin{go}
func (r RepoPin) Resolve(raw string) string {
    abs := raw
    if !filepath.IsAbs(abs) {
        abs = filepath.Join(r.Root, abs)
    }
    abs = filepath.Clean(abs)
    resolved := resolveExisting(abs)
    root := resolveExisting(r.Root)
    if resolved == root || strings.HasPrefix(resolved, root+string(filepath.Separator)) {
        return abs
    }
    return ""
}
\end{go}

\section{The primitive runs}

\file{internal/prim/fs.go}. \code{fs-edit} replaces exactly one occurrence.
Its \code{Preview} renders the diff the confirm prompt shows; its handler
re-resolves the path (defense in depth: the guard already denied an escape)
and refuses ambiguous or missing matches with messages written for the
model to act on.

\begin{go}
func applyEdit(content, old, replacement string) (string, error) {
    if old == "" {
        return "", errors.New("old text must not be empty")
    }
    switch n := strings.Count(content, old); n {
    case 0:
        return "", errors.New("old text not found; it must match exactly, including whitespace")
    case 1:
        return strings.Replace(content, old, replacement, 1), nil
    default:
        return "", fmt.Errorf("old text occurs %d times; include more context so it is unique", n)
    }
}
\end{go}

\section{The result is recorded, and shown}

Back in \file{loop.go}, \code{runTool} records \code{tool.call} before and
\code{tool.result} after. Only context cancellation is an error here;
everything else\,---\,a denial, a missing file, a bad argument\,---\,is a
\emph{result} the model reads:

\begin{go}
res := log.ToolResult{ID: tc.ID, Name: tc.Name}
value, err := l.apply(ctx, tc)
if err != nil {
    if ctx.Err() != nil {
        return ctx.Err()
    }
    res.OK, res.Error = false, err.Error()
} else {
    res.OK, res.Value = true, value
}
_, err = l.host.Record(log.KindToolResult, res)
\end{go}

And \file{internal/log/log.go} writes it: marshal, filter (the future
redactor), one line, fsync.

\begin{go}
line, err := json.Marshal(ev)
line = append(line, '\n')
if _, err := w.w.Write(line); err != nil {
    return Event{}, fmt.Errorf("log: write: %w", err)
}
if w.syncer != nil {
    if err := w.syncer.Sync(); err != nil {
        return Event{}, fmt.Errorf("log: fsync: %w", err)
    }
}
w.seq = ev.Seq
\end{go}

The front-end, meanwhile, has been receiving each of these as an event and
printing a tool card, a guard flag, or a dimmed diff. That is the whole
trip.

\chapter{The log as truth}

A session log lives under \file{\textasciitilde/.local/share/nodwyn/sessions/}, one \file{<id>.jsonl} per session. The
kinds are fixed in \file{internal/log/log.go}; a reader rejects any other,
and rejects a gap in \code{seq}. Three lines from a real session:

\begin{txt}
{"seq":2,"ts":"2026-09-21T03:14:52Z","session":"20260921T031452Z-da2e39d8",
 "kind":"user.msg","payload":{"text":"Read internal/log/log.go and explain ..."}}
{"seq":5,"ts":"...","kind":"tool.call","payload":{"id":"call_x","name":"fs-read",
 "args":{"path":"internal/log/log.go"},"origin":"model"}}
{"seq":6,"ts":"...","kind":"guard.decision","payload":{"prim":"fs-read",
 "origin":"model","decision":"allow","args":"(\"internal/log/log.go\")"}}
\end{txt}

What is \emph{not} in the log is as telling as what is: \code{llm.req} records
how many messages and an estimate, not the messages\,---\,they are
reconstructible. Model thinking is stored on \code{llm.resp} and never sent
back. Pure builtins (\code{car}, \code{+}) are allowed silently; anything that
reaches the host, and every non-allow, is recorded.

Two properties come for free once the log is the truth. A crash between
write and fsync leaves a partial last line; \code{Open} repairs it and says
so. And ``resume'' (not built yet) is \code{Build} over an old file.

\chapter{Conventions}

\paragraph{Metadata is the contract.} A primitive declares \code{Doc},
\code{Params} (name, type, doc, \emph{role}), \code{Variadic}, \code{Egress},
\code{Irreversible}, \code{Host}, \code{Hidden} and an optional \code{Preview}.
The model's tool schema and the guard's judgement both come from this one
declaration and cannot drift. Roles\,---\,\code{path}, \code{argv},
\code{payload}, later \code{dest} and \code{credential}\,---\,are what checkers
read. A checker that switches on a primitive's name is a bug.

\paragraph{Origins.} A context carries who initiated the evaluation:
\code{user} (REPL, \code{:eval}, \file{init.scm}) or \code{model} (a tool
call). Every guard record has it; L2 will use it to restrict tag stripping.

\paragraph{Errors.} \code{*scheme.Error} is the Scheme-level error, with an
optional \code{Cause} for host errors (a \code{*Denied}, a cancelled
context). Runtime failures read \code{who: what}. A tool's failure is a
result for the model; only cancellation is an error for the loop.

\paragraph{Two binaries.} \file{bin/nodwyn} is what you are editing;
\code{nodwyn-stable} is what edits it. \code{make promote} is the only bridge.

\paragraph{One ingress, one exit.} There is one function that writes the
log, one that appends to the model's context (the window), one that calls a
primitive's handler. Lints and tests hold each. When a phase needs a second
one, the answer is an ADR, not a shortcut.

\part{\nibb{} Scheme}

\chapter{Why a Scheme, and which one}

\nibb{} embeds a small Scheme for \emph{orchestration only}. Scheme code
decides what to do; every effect\,---\,a file, a process, a model call\,---\,is a
Go primitive, and every primitive application crosses the guard. Scheme
never touches the network or the filesystem on its own. That gives the
harness one place to look (\adr{3}, invariant~4) and gives the user a way to
extend it without touching Go: a tool is a documented Scheme procedure
registered into the primitive table.

The language is deliberately small. It is specified in
\file{docs/scheme.md} and the interpreter implements that document and
nothing else; the conformance suite has a snippet for every form and
builtin, and a builtin that is not documented fails the suite. What follows
is a tutorial over that reference.

\begin{aside}
\textbf{On P1b.} The forms \code{letrec}, \code{case}, \code{do},
\code{when}, \code{unless}, \code{guard}, \code{with-output-to-string},
vectors and block comments arrived in Phase 1b, the first phase after the
bootstrap. Their sections are marked \textbf{[P1b]} so a reader of the
bootstrap code knows where the kernel grew.
\end{aside}

\chapter{First steps}

\section{The REPL}

\begin{sh}
$ nibb repl
nibb> (+ 1 2)
3
nibb> (define (square x) (* x x))
nibb> (map square (iota 5))
(0 1 4 9 16)
nibb> (car 1)
error: car: argument 1: expected pair, got integer 1
\end{sh}

Unbalanced input continues on the next line with a \code{...>} prompt. A
value of \emph{void}\,---\,the result of \code{define}, \code{set!},
\code{display}\,---\,is not printed. \code{nibb eval -e EXPR} evaluates one
expression; \code{nibb eval FILE} loads files. Inside a session,
\code{:eval EXPR} evaluates with \emph{user} origin, and the guard sees it
like anything else.

\section{Values}

\begin{center}
\small
\begin{tabular}{@{}llll@{}}
\toprule
Type & Written & Predicate & Notes \\
\midrule
empty list & \code{()} & \code{null?} & quote it in code: \code{'()} \\
boolean & \code{\#t} \code{\#f} & \code{boolean?} & only \code{\#f} is false \\
integer & \code{42} & \code{integer?} & int64; wraps on overflow \\
float & \code{2.5} \code{1.0} & \code{float?} & float64; always printed with a dot or exponent \\
string & \code{"hi"} & \code{string?} & immutable UTF-8; lengths and indices count runes \\
symbol & \code{foo} & \code{symbol?} & case-sensitive \\
pair / list & \code{(1 2 3)} \code{(a . b)} & \code{pair?} \code{list?} & no mutation \\
vector [P1b] & \code{\#(1 2)} & \code{vector?} & \code{vector-set!} mutates \\
procedure & \code{\#<procedure f>} & \code{procedure?} & closures and primitives alike \\
handle & \code{\#<cred:class:id>} & \code{handle?} & opaque; no builtin turns it into text \\
void & \code{\#<void>} & --- & the unspecified value \\
\bottomrule
\end{tabular}
\end{center}

\code{number?} is true for both numeric kinds. Two integers divide exactly
when the result is integral, otherwise the result is a float:

\begin{scm}
(/ 6 3)      ; 2
(/ 7 2)      ; 3.5
(/ 1.0 4)    ; 0.25
(expt 2 10)  ; 1024
(sqrt 16)    ; 4.0
(exact 2.9)  ; 2   (truncates toward zero)
(inexact 2)  ; 2.0
\end{scm}

\section{Printing}

\code{display} shows a value for humans (strings without quotes);
\code{write} shows it for machines (strings quoted and escaped so that the
reader can read them back). Both go to the interpreter's output writer,
which the host provides: stdout in the REPL, an \emph{output} event in a
session. \code{format} builds a string with \code{\textasciitilde a}
(display), \code{\textasciitilde s} (write), \code{\textasciitilde\%}
(newline) and \code{\textasciitilde\textasciitilde}:

\begin{scm}
(format "~a and ~s~%" "x" "y")   ; "x and \"y\"\n"
\end{scm}

\chapter{Definitions and procedures}

\section{define and lambda}

\begin{scm}
(define pi 3.14159)
(define (area r) (* pi r r))
(define (rest-args a . more) more)    ; dotted parameter: the rest as a list
(define variadic (lambda args args))  ; every argument as a list

(area 2)               ; 12.56636
(rest-args 1 2 3)      ; (2 3)
(variadic 1 2)         ; (1 2)
\end{scm}

A \code{define} inside a procedure body binds in that call's environment.
When \code{define} names a fresh anonymous closure, the closure takes that
name for printing and for error messages: \code{(define f (lambda (a) a))
(f)} reports \code{f: expected 1 argument(s), got 0}.

\section{Closures}

Procedures close over their defining environment. \code{set!} assigns an
\emph{existing} binding; assigning an unbound name is an error.

\begin{scm}
(define (counter)
  (let ((n 0))
    (lambda () (set! n (+ n 1)) n)))
(define c (counter))
(c) (c) (c)   ; 3
\end{scm}

\section{let, let* and named let}

\code{let} binds in parallel, \code{let*} sequentially. \emph{Named let}
is the idiomatic loop: the name is bound to the loop procedure inside the
body.

\begin{scm}
(let ((a 1) (b 2)) (+ a b))            ; 3
(let* ((a 1) (b (+ a 1))) b)           ; 2
(let loop ((i 0) (acc '()))
  (if (= i 3)
      (reverse acc)
      (loop (+ i 1) (cons i acc))))    ; (0 1 2)
\end{scm}

\section{Tail calls}

The last expression of a body, of \code{begin}, of a \code{let} body, of a
\code{cond} clause, of \code{and}/\code{or}, and either branch of \code{if}
is in \emph{tail position}. A call there does not grow the stack: a loop
written as a tail-recursive procedure runs in constant space, for a
million iterations or a billion. The test suite proves this with a
200,000-iteration loop under an 8\,MB stack cap in nine positions.

\begin{scm}
(define (count-down n)
  (if (= n 0) 'done (count-down (- n 1))))
(count-down 1000000)    ; done
\end{scm}

Mutual recursion works the same way:

\begin{scm}
(define (even2? n) (if (= n 0) #t (odd2? (- n 1))))
(define (odd2? n)  (if (= n 0) #f (even2? (- n 1))))
(even2? 1000000)        ; #t
\end{scm}

\section{letrec [P1b]}

\code{letrec} binds names that may refer to each other, for local mutual
recursion. Initialisers should be lambdas; a plain value must not depend on
a sibling that has not been evaluated yet.

\begin{scm}
(letrec ((ev? (lambda (n) (if (= n 0) #t (od? (- n 1)))))
         (od? (lambda (n) (if (= n 0) #f (ev? (- n 1))))))
  (ev? 10))   ; #t
\end{scm}

\chapter{Conditionals}

\begin{scm}
(if test then)                 ; missing else yields void
(if test then else)
(cond (test body ...)
      (test => proc)           ; pass the test's value to proc
      (else body ...))
(and a b c)                    ; short-circuit; (and) is #t
(or a b c)                     ; short-circuit; (or) is #f
\end{scm}

\code{cond} with \code{=>} is the idiom for ``look it up, and if found, use
it'':

\begin{scm}
(cond ((assq 'b '((a 1) (b 2))) => cadr)
      (else 0))    ; 2
\end{scm}

\textbf{[P1b]} \code{when} and \code{unless} evaluate a body for effect;
\code{case} dispatches on a key against lists of datums:

\begin{scm}
(when (> x 0) (display "positive") (newline))
(unless done? (display "still going"))
(case 2 ((1 2 3) 'small) ((4 5) 'medium) (else 'large))   ; small
(case 'Foo ((foo) 'lower) (else 'exact))                 ; exact -- symbols compare by name
\end{scm}

\chapter{Lists and association lists}

\section{Pairs and lists}

\begin{scm}
(cons 1 2)                 ; (1 . 2)
(list 1 2 3)               ; (1 2 3)
(car '(1 2 3))             ; 1
(cdr '(1 2 3))             ; (2 3)
(cadr '(1 2 3))            ; 2      also caar cdar cddr caddr cdddr
(length '(1 2 3))          ; 3
(append '(1) '(2 3) '())   ; (1 2 3)
(reverse '(1 2 3))         ; (3 2 1)
(list-ref '(a b c) 1)      ; b
(list-tail '(a b c) 2)     ; (c)
(iota 4)                   ; (0 1 2 3)
\end{scm}

Lists are not mutated in \nibb{} Scheme: there is no \code{set-car!}. Build
new ones.

\section{Higher-order}

\begin{scm}
(map (lambda (x) (* x x)) '(1 2 3))        ; (1 4 9)
(map + '(1 2) '(10 20 30))                 ; (11 22)   stops at the shortest
(for-each display '(1 2))                  ; prints 12, yields void
(filter odd? '(1 2 3 4))                   ; (1 3)
(fold + 0 '(1 2 3))                        ; 6         (f acc x), left to right
(fold cons '() '(1 2))                     ; ((() . 1) . 2)
(sort '(3 1 2) <)                          ; (1 2 3)   stable
(apply + 1 2 '(3 4))                       ; 10        last argument is a list
\end{scm}

\section{Membership and association lists}

\code{memq} and \code{assq} compare with \code{eq?}; \code{member} and
\code{assoc} with \code{equal?}. \code{eq?} is identity for pairs and
procedures and \emph{value} for everything else, including strings\,---\,so
\code{(eq? "a" "a")} is true, unlike in most Schemes. \code{equal?} is
structural.

\begin{scm}
(memq 'b '(a b c))                   ; (b c)
(member '(1) '((1) 2))               ; ((1) 2)
(assq 'b '((a . 1) (b . 2)))         ; (b . 2)
(assoc "k" '(("k" . 1)))             ; ("k" . 1)
\end{scm}

Association lists are \nibb's record type\,---\,JSON objects become alists,
tool metadata is an alist, \code{shell-exec} returns one. The prelude adds
\code{assoc-ref}:

\begin{scm}
(define r '(("status" . 0) ("stdout" . "ok\n") ("stderr" . "")))
(assoc-ref r "status")          ; 0
(assoc-ref r "missing")         ; #<void>
(assoc-ref r "missing" 'none)   ; none
\end{scm}

\chapter{Strings, JSON and quasiquote}

\section{Strings}

\begin{scm}
(string-length "héllo")                  ; 5   (runes)
(string-append "a" "b" "c")              ; "abc"
(substring "héllo" 1 3)                  ; "él"
(string-upcase "aé")                     ; "AÉ"
(string-trim "  x \n")                   ; "x"
(string-split "a,b,,c" ",")              ; ("a" "b" "" "c")   literal separator
(string-join '("a" "b") "-")             ; "a-b"
(string-contains? "hello" "ell")         ; #t
(string-prefix? "he" "hello")            ; #t
(string-replace "a-b-c" "-" "+")         ; "a+b+c"
(string->number "42")                    ; 42       #f if not a number
(number->string 1.5)                     ; "1.5"
(string->symbol "ab") (symbol->string 'ab)
\end{scm}

There is no character type; one-character strings serve.

\section{JSON}

\code{json-parse} and \code{json-emit} are the bridge to the model and to
tools. Objects become alists with string keys in document order, arrays
lists, \code{null} the symbol \code{null}:

\begin{scm}
(json-parse "{\"a\":1,\"b\":[true,null,2.5]}")
; (("a" . 1) ("b" #t null 2.5))

(json-emit '(("name" . "nibb") ("tags" "go" "scheme")))
; "{\"name\":\"nibb\",\"tags\":[\"go\",\"scheme\"]}"
\end{scm}

\code{json-emit} treats a proper list whose elements are all pairs with
string cars as an object, and any other proper list as an array. So
\code{'()} emits \code{[]}, and there is no way to emit \code{\{\}} in v0.
This same mapping is how a model's tool-call arguments become Scheme values.

\section{Quasiquote}

One level, with \code{,} to evaluate and \code{,@} to splice:

\begin{scm}
`(1 ,(+ 1 1) ,@(list 3 4))   ; (1 2 3 4)
\end{scm}

\chapter{Errors, output capture and vectors}

\section{Raising}

\code{(error message irritant ...)} raises an error whose text is the
message followed by the written irritants. Runtime failures use the same
kind of error with a \code{who: what} message. An error that is not caught
unwinds to the host: the REPL prints it; in a session, the model receives it
as the tool result. A primitive the guard refuses raises
\code{denied: <reason>}.

\begin{scm}
(error "boom" 1 "two")    ; error: boom 1 "two"
\end{scm}

\section{Catching [P1b]}

\code{guard} has the R7RS shape: a variable, \code{cond}-style clauses, and
a body. If the body raises, the variable is bound to the \emph{condition}\,---\,an
alist with the message and the irritants\,---\,and the clauses are tried in
order; the first true clause's value is the value of the form. No match
re-raises the original error.

\begin{scm}
(guard (e ((string-contains? (assoc-ref e 'message) "boom") 'caught)
          (else (assoc-ref e 'irritants)))
  (error "boom" 1 2))                        ; caught

(guard (e (#t (assoc-ref e 'message)))
  (car 1))                                   ; "car: argument 1: expected pair, got integer 1"

(guard (e ((assoc-ref e 'message) => string-length))
  (error "four"))                            ; 4
\end{scm}

Two things are never catchable, whatever the clauses say: a guard
\emph{denial} and a \emph{cancellation}. A tool cannot hide that it was
refused, and an abort always unwinds. The test suite checks both.

\section{Capturing output [P1b]}

\begin{scm}
(with-output-to-string (display "hi") (newline))   ; "hi\n"
\end{scm}

\section{Vectors [P1b]}

\begin{scm}
(vector 1 "a" #t)                 ; #(1 "a" #t)
(vector-ref #(10 20 30) 1)        ; 20
(let ((v (vector 1 2 3)))
  (vector-set! v 1 99)
  (vector->list v))               ; (1 99 3)
(equal? #(1 2) #(1 2))            ; #t   eq? compares identity
\end{scm}

\chapter{Primitives and the guard, from Scheme}

Every builtin, every host primitive and every registered tool is an entry in
the \emph{primitive table}. You can inspect it:

\begin{scm}
(length (primitives))                        ; ~130
(primitive-info "fs-edit")
; ((doc . "Edit a text file ...") (params . (("path" "string" "...") ...))
;  (variadic . #f) (egress . #f) (irreversible . #t) (host . #t) (origin . "builtin"))
\end{scm}

The host primitives in the bootstrap are \code{fs-read}, \code{fs-write},
\code{fs-edit}, \code{fs-list} and \code{shell-exec}. Calling one from
Scheme is calling it: the guard runs, and if it says no you get an error.

\begin{scm}
(fs-read "README.md")                   ; the contents
(fs-list ".")                           ; ("Makefile" "cmd/" ...)   dirs end in /
(fs-read ".env")                        ; error: fs-read: denied: path .env matches the credential-file denylist (.env)
(shell-exec "go" '("version"))          ; (("status" . 0) ("stdout" . "go version ...") ("stderr" . ""))
\end{scm}

\code{shell-exec} takes the program and a \emph{list} of arguments\,---\,there is
no shell, no globbing, no \code{\$VAR}. The child runs in the repository
root with an allowlisted environment; its output is capped; cancellation
kills the whole process group.

\chapter{Writing tools}

\section{register-tool}

A tool is a Scheme procedure plus metadata, registered into the primitive
table. From then on it is indistinguishable from a builtin at the dispatch
site, and the model sees it in its tool list.

\begin{scm}
(register-tool "go-test"
  (lambda (pkg) (run "go" (list "test" "-race" pkg)))
  '((doc . "Run the Go tests with the race detector for a package pattern.")
    (params . (("pkg" "string" "package pattern, e.g. ./internal/log or ./...")))))
\end{scm}

The metadata alist:

\begin{center}
\small
\begin{tabular}{@{}lp{9.5cm}@{}}
\toprule
Key & Meaning \\
\midrule
\code{doc} & one line; required; shown to the model as the tool description \\
\code{params} & list of \code{("name" "type" "doc")}; required; type is one of \code{string integer number boolean list alist vector any} \\
\code{variadic} & extra arguments allowed after \code{params}; default \code{\#f}. Prefer a list parameter: the model's schema cannot express rest arguments \\
\code{egress}, \code{irreversible} & default \code{\#f}; declare them honestly, the guard reads them \\
\bottomrule
\end{tabular}
\end{center}

Rules: the name must be new\,---\,builtins cannot be shadowed and a tool
cannot be re-registered, because the tool list must stay stable for the
whole session (\adr{9}). A registered tool's own primitive calls are guarded
in turn: a wrapper over \code{shell-exec} does not make \code{shell-exec}
less irreversible.

\section{The dev tools, read closely}

\file{dev/go.scm} is the pattern. A helper renders a \code{shell-exec}
result the way a model reads best\,---\,exit status, then the output:

\begin{scm}
(define (run cmd args)
  (let ((r (shell-exec cmd args)))
    (format "exit ~a~%~a~a"
            (assoc-ref r "status")
            (assoc-ref r "stdout")
            (assoc-ref r "stderr"))))
\end{scm}

\file{dev/repo.scm} shows a tool that needs no shell at all: find an ADR by
number by listing a directory and reading a file.

\begin{scm}
(define (adr-file n)
  (let ((prefix (format "~a-" (pad4 n))))
    (let loop ((names (fs-list "docs/adr")))
      (cond ((null? names) #f)
            ((string-prefix? prefix (car names))
             (string-append "docs/adr/" (car names)))
            (else (loop (cdr names)))))))

(register-tool "read-adr"
  (lambda (n)
    (let ((f (adr-file n)))
      (if f (fs-read f) (error "no such ADR" n))))
  '((doc . "Read architecture decision record number n from docs/adr.")
    (params . (("n" "integer" "the ADR number, e.g. 10")))))
\end{scm}

\section{init.scm and allow.toml}

\file{<repo>/.nibb/init.scm} is loaded at start, after the prelude, with
user origin and with confirms denied (nobody is listening yet), so it may
register tools and must not act. \code{load} resolves relative paths against
the repository root:

\begin{scm}
(load "dev/go.scm")
(load "dev/git.scm")
(load "dev/repo.scm")
\end{scm}

Pair each command-wrapping tool with a \emph{standing approval} in
\file{.nibb/allow.toml}, or every call prompts:

\begin{txt}
commands = ["go test", "go build", "git status", "git diff", "git grep"]
writes   = ["internal/", "docs/"]
\end{txt}

\chapter{Reference}

\section{Special forms}

\begin{center}
\small
\begin{tabular}{@{}p{7.6cm}p{6.2cm}@{}}
\toprule
Form & Meaning \\
\midrule
\code{(define name expr)} & bind in the current environment \\
\code{(define (name p ... [. rest]) body ...)} & named lambda \\
\code{(lambda (p ... [. rest]) body ...)}, \code{(lambda args body ...)} & closure \\
\code{(if test then [else])} & missing else yields void \\
\code{(let ((v e) ...) body ...)} & parallel bindings \\
\code{(let name ((v e) ...) body ...)} & named let \\
\code{(let* ((v e) ...) body ...)} & sequential bindings \\
\code{(letrec ((v e) ...) body ...)} [P1b] & mutually recursive bindings \\
\code{(begin e ...)} & sequence \\
\code{(set! name expr)} & assign an existing binding \\
\code{(quote d)}, \code{'d} & the datum \\
\code{(quasiquote d)}, \code{`d} with \code{,e} \code{,@l} & one level \\
\code{(cond (test body ...) ...} \code{(test => proc) (else ...))} & first true clause \\
\code{(case key ((d ...) body ...) ... (else ...))} [P1b] & dispatch on datums \\
\code{(and e ...)}, \code{(or e ...)} & short-circuit \\
\code{(when test e ...)}, \code{(unless test e ...)} [P1b] & conditional effect \\
\code{(do ((v init [step]) ...)} \code{(test [e ...]) cmd ...)} [P1b] & loop, constant stack \\
\code{(guard (var clause ...) body ...)} [P1b] & catch errors; denials and cancellation excepted \\
\code{(with-output-to-string e ...)} [P1b] & capture output \\
\bottomrule
\end{tabular}
\end{center}

Special-form names are reserved: binding one as a variable does not change
how the form is parsed.

\section{Builtins by group}

\begin{txt}
Numbers      + - * / = < > <= >= quotient remainder modulo abs min max
             floor ceiling round truncate exact inexact expt sqrt
             number? integer? float? zero? positive? negative? even? odd?
             number->string string->number
Equivalence  not boolean? eq? eqv? equal?
Lists        cons car cdr caar cadr cdar cddr caddr cdddr list list? pair?
             null? length append reverse list-ref list-tail iota memq member
             assq assoc map for-each filter fold sort apply
Vectors      vector vector? vector-length vector-ref vector-set!
[P1b]        vector->list list->vector
Strings      string? string-length string-append substring string=? string<?
             string-upcase string-downcase string-trim string-split string-join
             string-contains? string-prefix? string-suffix? string-replace
             string->symbol symbol->string symbol? format
Output       display write newline
Control      procedure? apply void error
JSON         json-parse json-emit
Meta         primitives primitive-info register-tool load handle?
Prelude      assoc-ref identity compose string-empty?        (Scheme)
Host         fs-read fs-write fs-edit fs-list shell-exec     (bootstrap)
\end{txt}

\section{Not in \nibb{} Scheme}

\code{set-car!}/\code{set-cdr!} and string mutation; nested quasiquote;
hygienic macros; \code{call/cc}; \code{dynamic-wind}; an \code{eval}
primitive (the kernel surface is not exposed to programs); characters;
exact rationals and bignums; source positions in runtime errors.

\part{Using \nibb}

\chapter{Install and configure}

\begin{sh}
make build                        # bin/nibb from the working tree
make promote TAG=v0.1-bootstrap   # accept + selfcheck, then ~/.local/bin/nibb-stable
\end{sh}

Use \code{nibb-stable} for real work and \file{bin/nibb} only to test a
change to \nibb{} itself. Requires Go~1.27, \code{staticcheck}, and a
reachable control plane.

\section{config.toml}

\file{\textasciitilde/.config/nibb/config.toml}; a missing file means these
defaults:

\begin{txt}
[control]
provider   = "ollama"
url        = "http://workshop.example:11434"
model      = "qwen3.8-27b-96k:latest"
num_ctx    = 98304
keep_alive = "30m"

[window]
high_water        = 61440  # trim above this (tokens; validated against num_ctx)
low_water         = 30720  # trim down to this, then hold
max_result_tokens = 3072   # larger tool results are shown head+tail
output_reserve    = 8192

[policy]
unbound        = "deny"    # or "confirm" (P3)
allow_commands = []        # standing approvals, global
allow_writes   = []

[loop]
request_timeout = "30m"
shell_timeout   = "10m"

[budget]
tool_calls_soft = 50
tool_calls_hard = 500
tokens_soft     = 1000000
tokens_hard     = 10000000
actions_soft    = 25
actions_hard    = 200
wall_soft       = "20m"
wall_hard       = "4h"
\end{txt}

\code{nibb paths} prints where everything lives. Session logs are under
\file{\textasciitilde/.local/share/nibb/sessions/}.

\chapter{A session}

\begin{sh}
$ cd ~/src/your-repo
$ nibb-stable
control plane: http://workshop...:11434 qwen3.8-27b-96k:latest (ollama 0.34.2)
log: /Users/you/.local/share/nibb/sessions/20260921T143846Z-beba937b.jsonl
nibb session 20260921T143846Z-beba937b
repo  /Users/you/src/your-repo
tools fs-edit fs-list fs-read fs-write shell-exec
:h for commands

you> read main.go and explain the startup sequence
\end{sh}

Commands at the prompt:

\begin{center}
\begin{tabular}{@{}ll@{}}
\toprule
\code{:q} & quit \\
\code{:log [n]} & last $n$ log events \\
\code{:tools} & tools offered to the model \\
\code{:allow} & standing command and write approvals in effect \\
\code{:eval EXPR} & evaluate Scheme yourself (user origin) \\
\code{:thinking} & toggle streaming the model's thinking \\
Ctrl-C & abort the running turn; at the prompt, exit \\
\bottomrule
\end{tabular}
\end{center}

While a turn runs the front-end prints streamed text, a card per tool call
(\code{> fs-read \{"path":...\}} then \code{< fs-read ok (10398 bytes, 5 ms)}),
guard flags, and a usage line per model request.

\chapter{Confirms}

Every write and every command stops for you, with a reason and, for
writes, a diff:

\begin{txt}
--- confirm: fs-edit ("internal/x.go" "old..." "new...")
  * fs-edit is irreversible
  * changes internal/x.go
  --- internal/x.go
  +++ internal/x.go
  @@ -12,3 +12,3 @@
  ...
allow? [y/N/a=always under internal//t=all writes this turn]
\end{txt}

\begin{description}[leftmargin=2.2em]
\item[\code{y}] once.
\item[\code{N}] (or anything else) no. The model is told
  \code{denied: not confirmed by the user} and moves on.
\item[\code{a}] yes, and a \emph{standing approval for this session}: the
  command prefix (\code{go test}, \code{rg}) or the write's directory
  (\code{internal/scheme/}). Not persisted.
\item[\code{t}] writes only: yes, and every further write \emph{until this
  turn ends}. For a task you intend to review with \code{git diff}
  afterwards. Each auto-approved write's diff is logged and shown dimmed.
\end{description}

Read the reason lines. ``command line names a path outside the repository''
and ``uses the network tool curl'' are the guard telling you what a
\code{y} would do.

\section{Standing approvals}

\begin{txt}
# ~/.config/nibb/config.toml           (global)
[policy]
allow_commands = ["go test", "go build", "git status", "git diff", "rg"]
allow_writes   = []

# <repo>/.nibb/allow.toml              (per repository, tracked)
commands = ["make lint", "make test"]
writes   = ["internal/", "docs/"]
\end{txt}

Prefixes match argv tokens: \code{"go test"} matches \code{go test ./...}
and \code{/usr/local/go/bin/go test}, not \code{go build}. Denies still
apply\,---\,only the confirm is waived\,---\,and every use is logged with the
prefix as reason. \code{.nibb/}, \code{.github/}, \code{.gitea/},
\code{scripts/} and \code{Makefile} always prompt, whatever the list says.
Keep the lists read-only and test-only for commands; git is the undo for
writes.

\chapter{Building \nibb{} with \nibb}

\begin{enumerate}[leftmargin=2em]
\item \code{nibb-stable} in this repository; \code{:tools} lists
  \code{go-test}, \code{git-diff}, \code{grep-repo}, \code{read-adr}, \ldots
\item Give it small, explicit tasks: name the file, the pattern to follow,
  the test to add, and ask it to run \code{go-test} on the package. A
  paragraph works; a page does not.
\item Review every diff at the confirm, or \code{t} and review with
  \code{git diff} at the end. An \code{fs-write} of a whole existing file
  deserves suspicion; an \code{fs-edit} rarely does.
\item When it gets in the way, hand-edit, and add one line to
  \file{docs/friction.md}.
\item \code{make accept} before committing. You commit; there is no commit
  tool.
\item \code{make promote TAG=v0.N-phase} when a phase's acceptance is green.
  Promotion refuses a dirty tree, runs the suite and the self-check, swaps
  \code{nibb-stable} (keeping \code{.prev}) and tags.
\end{enumerate}

\chapter{When something is wrong}

\begin{center}
\small
\begin{tabular}{@{}p{5.6cm}p{8.6cm}@{}}
\toprule
Symptom & Look at \\
\midrule
\code{control plane ... unreachable} & tailscale up? \code{curl http://workshop...:11434/api/version} \\
first turn takes 10--60\,s & another client holds workshop's single slot, or the prefix is cold; later turns are sub-second \\
the model loops on a denied call & it was told; say so in the next message; the hard stop is 50 tool calls \\
\code{hard stop: too many tool calls} & the turn ended; the log has everything; rephrase smaller \\
a tool call errored & \code{:log 20}: the \code{tool.result} line has the exact error the model saw \\
a file changed without a prompt & \code{:allow}: a standing approval covered it; the \code{guard.decision} line names which and carries the diff \\
a session died & \code{nibb replay <id>}: everything up to the last fsync is there \\
\bottomrule
\end{tabular}
\end{center}

\appendix
\chapter{Event kinds}

The log's kinds are fixed in \file{internal/log/log.go}; a reader rejects any
other, and rejects a gap in \code{seq}.

{\small
\begin{longtable}{@{}L{4.8cm}L{10.2cm}@{}}
\toprule
Kind & What it records \\
\midrule
\endhead
\code{session.start} & version, root, the control model, how many tools, the window recipe \\
\code{session.resume} & a new process continuing the same conversation \\
\code{session.end} & clean shutdown; absent after a crash \\
\code{user.msg} & text you sent, redacted \\
\code{harness.msg} & text the loop's strategy added for the model \\
\code{attach} & an attached file: name, type, hash, extracted and redacted text; for an image, its size, token estimate and confirmed host, never its pixels \\
\code{llm.req} & a request to a model: which model, what was elided, settings (the reasoning effort included), a digest, and the attached images it did not carry, with why \\
\code{llm.resp} & the model's reply: content, thinking, tool calls, stop reason, usage (tokens; cached tokens and cost when known) \\
\code{tool.call} & a tool application about to be evaluated, with its origin \\
\code{tool.result} & its result or error, redacted \\
\code{guard.decision} & primitive, origin, verdict, checker, reason, answer and who answered \\
\code{egress.check} & a binding lookup for a keyed request: class, host, method, an OAuth2 credential's scopes, outcome \\
\code{gate.plan} & the plan gate's outcome for a message \\
\code{plan.step} & the model marked a step of the plan done or skipped, with its note \\
\code{jev.req} & a call to the judge: which call, its origin, content level, size, payload hash, redactions, answer, the model that answered and its tokens (priced with \code{[judge]} prices, else the published ones) \\
\code{cred.register} & a credential stored: class and handle id only \\
\code{cred.refresh} & an OAuth2 credential refreshed before a use: class, handle id, token endpoint, expiry or error; never a token \\
\code{mcp.egress} & a contained MCP server's connection through its egress proxy, allowed or refused: server, method, host, port, why; past a limit, how many were not recorded \\
\code{redact.event} & a \code{[REDACTED class]} substitution and where it happened \\
\code{binding.change} & a binding row requested, applied or refused \\
\code{approval.change} & a standing approval granted explicitly \\
\code{turn.checkpoint} & a budget checkpoint and its answer \\
\code{turn.end} & the durable end of a turn \\
\code{run.state}, \code{run.turn}, \code{run.checkpoint} & a run's progress: its limits and state, whether it is unattended, auto-approved or judged, its turns, and each checkpoint's answer \\
\code{context.narrow}, \code{context.widen}, \code{context.fold}, \code{result.fold} & changes to what later windows show \\
\code{context.compact} & a compaction: what set it off, the context's size before and after, its handover note \\
\code{summary}, \code{summary.drop} & a revocable summary over a span, and its revocation \\
\code{file.version} & a file's previous content, kept in \file{.nodwyn/versions/} before a write: its path, size and hash \\
\code{doc.upload} & a document sent to Google Drive: file, size, hash, destination, the document's id or why nothing was sent; never content or a token \\
\code{setting.change} & a setting a person chose at a prompt, such as how uploads to Google Docs are handled: who, the value, whether \file{config.toml} took it \\
\code{session.title}, \code{session.archive} & the name a person gave the task, and whether it is archived (\adr{39}): who, and the name or the state; the latest of each counts \\
\code{session.model} & the model and reasoning effort the conversation talks to, chosen as it opened or changed (\adr{40}, \adr{41}): from, to, where it is reached, who chose it, and the window's recipe sized for it; one chosen while a turn ran is first recorded as queued, without a recipe; the latest one counts \\
\code{worktree.create}, \code{worktree.restore}, \code{worktree.remove} & the task's own git worktree (\adr{34}): made, with its repository, path, branch, base and the uncommitted files it does not have; made again by a resume; removed on close with its branch, when nothing in them was unlanded \\
\code{eval.req}, \code{eval.result} & Scheme you evaluated, and its redacted result \\
\code{note} & free-form: hard stops, standing approvals, turn trust \\
\bottomrule
\end{longtable}}

\begin{sloppypar}
Ephemeral events reach a front-end but are never logged:
\code{stream.text}, \code{stream.thinking}, \code{prompt.confirm},
\code{prompt.secret}, \code{prompt.checkpoint}, \code{prompt.closed},
\code{status}, \code{error}, \code{output}, \code{policy.result},
\code{inspection.result}.
\end{sloppypar}

\chapter{Decision records}

{\small
\begin{longtable}{@{}lp{9.4cm}l@{}}
\toprule
ADR & Title & Status \\
\midrule
\endhead
1 & Value representation: one tagged struct with a reserved taint field & accepted \\
2 & Log format and event kinds & accepted \\
3 & Model tool calls are Scheme applications & accepted \\
4 & Primitive metadata & accepted \\
5 & Keychain crypto & accepted \\
6 & Unbound egress defaults to deny & accepted \\
7 & Scheme scope: written from scratch, documented first & accepted \\
8 & LLM client is a Provider interface, Ollama first & accepted \\
9 & Context and cache discipline & accepted \\
10 & Core/front-end boundary & accepted \\
11 & Terminal stack & accepted \\
12 & Two binaries, always & accepted \\
13 & Promotion (amended: no tags) & accepted \\
14 & \file{.nodwyn/init.scm} convention & accepted \\
15 & Friction log & accepted \\
16 & Fallback is allowed and honest & accepted \\
17 & Standing command approvals (amended: writes) & accepted \\
18 & Turn budgets and long-running turns & partly built \\
19 & Jev sees content, never secrets (amends the spec) & accepted \\
20 & Jev as egress primitives; \code{register-check} & built \\
21 & A real UI after v0: daemon, protocol, clients & partly built \\
22 & Cowork: workspaces, attachments, MCP servers & built \\
23 & Untrusted MCP servers: isolation, pinning, fail-closed & built \\
24 & OAuth2 credentials: acquisition, refresh, scopes & built \\
25 & Log-derived ingress redaction and buffered model output & accepted \\
26 & The agent loop is a Scheme procedure & accepted \\
27 & Configured models and narrowed contexts & accepted \\
28 & A scratch directory beside the repository & accepted \\
29 & Optional tool parameters & accepted \\
30 & One app, three ways to reach it: desktop, browser, phone & accepted \\
31 & The project is renamed \nodwyn{} & accepted \\
32 & The default loop works the way pil did & accepted \\
33 & Go 1.26 is the base & accepted \\
34 & Sessions in git worktrees & partly implemented \\
35 & A chat is the way to work, and it asks less (amends 17, 18, 30) & accepted \\
\bottomrule
\end{longtable}}

\chapter{Glossary}

\begin{description}[leftmargin=2.2em]
\item[approval] \code{[approval]} in \file{config.toml} (Settings
  \textrightarrow{} Approvals): whether writes in the folder go ahead
  (\code{writes}) and who answers a command first, the judge, you or the
  setting (\code{commands}), in a session a person answers.
\item[attachment] a file you brought into a conversation: copied into the
  workspace, its text extracted, redacted and capped, and shown to the model
  as data. An image is sent as an image, after a confirm of its own.
\item[binding] a user-owned row authorising a credential class for a
  destination, positions and methods; the egress authority.
\item[checker] one stage of the guard; reads a primitive's metadata and
  arguments, returns a verdict.
\item[control model] the model \code{(model-turn)} asks, configured as
  \code{[control]}.
\item[contained server] an MCP server \nodwyn{} runs in a container, from
  an image pinned by digest, with no network or only its egress proxy.
\item[effect] your declaration of what an MCP tool does: \code{read} or
  \code{write}; undeclared is treated as egress and irreversible.
\item[egress proxy] \code{nodwyn-egress}, the only way out of a contained
  server's network: it lets through the hosts listed and reports every
  attempt.
\item[handle] an opaque reference to a keychain entry; a Scheme value that
  can never become text.
\item[Jev] TypeSafe's external calibrated judge, one of the judge's two
  providers with OpenAI's Decisions API (\adr{46}), admitted at Seam A
  (plans against asks) and Seam C (egress situations), over redacted
  content, never a secret value.
\item[manifest] an MCP server's tool list; its hash is what you approve.
\item[model server] a provider's address and key, configured as
  \code{[servers.NAME]}; the models on it name it with \code{server}.
\item[mount] a folder outside a workspace's root that its sessions may
  read, or read and write.
\item[origin] who initiated an evaluation: the user, the model, or the
  harness (the loop's strategy).
\item[primitive] a Go handler in the table; the only way Scheme reaches the
  host.
\item[promotion] the gated copy of a candidate binary to \code{nodwyn-stable}.
\item[provenance tag] where a Scheme value came from (\code{external:file},
  \code{cred:openai-key}, \ldots); the union of a call's arguments flows into
  its result.
\item[security confirm] a confirm no standing approval, \code{[approval]}
  setting or judge can answer: it goes to a person, and a run nobody
  watches refuses it.
\item[standing approval] a command prefix or write path the user authorised
  ahead of time, and, by default, any write in the folder
  (\code{approval.writes}); converts an ordinary confirm to an allow, never
  a deny.
\item[strategy] the Scheme procedure that runs a turn, \code{agent-turn}.
\item[window] the model's context, rebuilt from the log on every request.
\item[workspace] a session's root plus its mounts.
\end{description}

\end{document}
